\begin{Solbook}{1}
 \par \noindent  \begin {listEX}[3] \item ĐKXĐ: $x \ne \pm 2$. \par $ \dfrac {3}{x+2} + \dfrac {x}{x-2} = \dfrac {x^2 - 1}{x^2 - 4}$ \\ $ 3(x - 2) + x(x + 2) = x^2 - 1$\\ $ 5x = 5 \\ x = 1 \text { (nhận)}$ \par Vậy phương trình đã cho có nghiệm $x = 1$. \par \item $\begin {cases} 5x - 3y = 13 \\ -2x + y = 6\end {cases} $\\ $\begin {cases} 5x - 3y = 13 \\ -6x + 3y = 18\end {cases}$ \\ $\begin {cases} -x = 31 \\ -6x + 3y = 18\end {cases}$ \\ $\begin {cases} x = -31 \\ y = -56\end {cases}$ \par Vậy hệ phương trình đã cho có nghiệm duy nhất\\ $(x; y) = (-31; -56)$. \par \item $7x + 2 < 3x - 2$ \\ $7x - 3x < -2 - 2$ \\ $4x < -4$ \\ $x < -1$ \par Vậy nghiệm của bất phương trình là $x < -1$. \end {listEX}  \hfill \qedEX 
\end{Solbook}
\begin{Solbook}{2}
 \par \noindent  \begin {enumerate} \par \item $2\sqrt {27} + 3\sqrt {12} - \sqrt {48} + \dfrac {3}{5}\sqrt {75}$ $= 6\sqrt {3} + 6\sqrt {3} - 4\sqrt {3} + 3\sqrt {3} = 11\sqrt {3}$. \item $\sqrt {(x - 2)^2} + \sqrt {(x - 5)^2}$ $= |x - 2| + |x - 5| = (x - 2) + (5 - x) = 3$, với $2 < x < 5$. \end {enumerate}  \hfill \qedEX 
\end{Solbook}
\begin{Solbook}{3}
 \par \noindent Gọi $x,y$ (chuyến) lần lượt là số chuyến của ô tô thứ hai và ô tô thứ ba. ($x,y \in \mathbb {N}^*$) \\ Xe thứ nhất chở được: $1{,}5x$ (chuyến). \\ Tổng cộng hết 50 chuyến, ta có: $1{,}5x + x + y = 50$ hay $2,5x+y=50$ \quad (1) \\ Tổng khối lượng hàng là 118 tấn, ta có: $2. 1{,}5x + 2{,}5x + 3y = 118$ hay $5,5x+3y=118$ \quad (2) \\ Từ (1) và (2), ta có hệ phương trình: $\begin {cases} 2{,}5x + y = 50 \\ 5{,}5x + 3y = 118 \end {cases}$ \par Giải hệ phương trình ta được: $x=16, y=10$ (nhận) \par Vậy ô tô thứ nhất chở 24 chuyển, ô tô thứ hai chở 16 chuyến và ô tô thứ ba chở 10 chuyến. \hfill \qedEX 
\end{Solbook}
\begin{Solbook}{4}
 \par \noindent  \begin {enumerate} \item Xét $\triangle MNK$ vuông tại N, ta có: \par $\tan K = \dfrac {MN}{KN}$ suy ra $KN = \dfrac {1200}{\tan 35^\circ } \approx 1714\, (\text {m})$. \par Vậy khoảng cách $KN$ từ người đó đến phương thẳng đứng của trực thăng khoảng 1714 m. \par \item Đổi $6$ phút $= \dfrac {1}{10}$ giờ. \par Xét $\triangle MNI$ vuông tại N: \par $\cos M = \dfrac {MN}{MI}$ suy ra $MI = \dfrac {1200}{\cos 20^\circ } (\text {m})=\dfrac {1,2}{\cos 20^\circ }$ (km). \par Vận tốc trung bình khi đáp xuống: $\dfrac {1,2}{\cos 20^\circ }: \dfrac {1}{10} \approx 12{,}8\, (\text {km/h})$. \end {enumerate} \hfill \qedEX 
\end{Solbook}
\begin{Solbook}{5}
 \par \noindent  \begin {listEX}[2] \item Ta có $\begin {cases} AB = AC (\text {gt}) \\ OB = OC (=R) \end {cases}$ \par nên $OA$ là đường trung trực của $BC$. \\ Suy ra $OA \perp BC$. \\ Mà $OA \perp AE$ (vì $AE$ là tiếp tuyến của $(O)$) \\ Suy ra $AE \parallel BC$ \quad (1) \item Xét $\triangle {DAE}$ và $\triangle {DCB}$ có: \par $\widehat {ADE} = \widehat {CDB}$ (đối đỉnh), \par $DA = DC$ (gt), \par $\widehat {DAE} = \widehat {DCB}$ (so le trong và $AE \parallel BC$). \par Suy ra $\triangle DAE = \triangle DCB$ (g-c-g) \par Suy ra $AE = BC$ (2 cạnh tương ứng) \quad (2) \\ Từ (1) và (2) suy ra tứ giác $ABCE$ là hình bình hành. \item ! Ta có $I$ là trung điểm của $CF$ suy ra $OG \perp CF$. \immini { Mà $CF \parallel AB$ (vì $ABCE$ là hình bình hành) \par Nên $OG \perp AB$ \par Suy ra $\widehat {BGO} = 90^\circ - \widehat {ABC}$ \par Mặt khác: $\widehat {BAC} = 180^\circ - 2\widehat {ABC} = 2.(90^\circ - \widehat {ABC})$ \par Vậy $\widehat {BAC}=2\widehat {BGO}$ }{ \begin {tikzpicture}[font=\footnotesize ] \def \R {2} \coordinate (O) at (0,0); \coordinate (A) at (0,\R ); \coordinate (B) at (-120:2); \coordinate (C) at (-60:2); \coordinate (D) at ($(A)!0.5!(C)$); \draw (O) circle(\R ); \path [name path=ax] (A)--(4,2); \path [name path=bd] (B)--($(B)!2!(D)$); \path [name intersections={of=ax and bd,by=E}]; \path [name path=ce] (C)--(E); \path [name path=circleO] (O) circle(2); \path [name intersections={of=ce and circleO,by={F1,F2}}]; \coordinate (F) at (F1); \coordinate (I) at ($(C)!0.5!(F)$); \path [name path=bc] (B)--($(B)!5!(C)$); \path [name path=oi] (O)--($(O)!5!(I)$); \path [name intersections={of=bc and oi,by=G}]; \draw (C)--(A)--(B)--(G)--(O) (C)--(E)--(A)--(O) (B)--(E); \foreach \x /\g in{O/-90,A/90,B/-180,C/-60,D/-30,E/90,F/0,I/0,G/0} \draw [fill=white] (\x ) circle (1.5pt) +(\g :.3) node{$\x $}; \end {tikzpicture} } \end {listEX} \hfill \qedEX 
\end{Solbook}
\begin{Solbook}{6}
 \par \noindent  $(x+2)(x+3)(x-5)(x-6)=180$ \\ $\left [(x+2)(x-5)\right ]\left [(x+3)(x-6)\right ]=180\\$ $(x^2-3x-10)(x^2-3x-18)=180$\\ Đặt $t=x^2-3x-14$\\ Khi đó pt trở thành: \( \begin {aligned}[t] (t+4)(t-4)&=180\\ t^2-16&=180\\ t^2&=196\\ t&=\pm 14 \end {aligned} \) \begin {multicols}{2} TH1:\,$t=14$\\ $x^2-3x-14=14$\\ $x^2-3x-28=0$\\ $(x-7)(x+4)=0$\\ $x=7\, \text {hoặc}\, x=-4$ \par \columnbreak \par TH2:\,$t=-14$\\ $x^2-3x-14=-14$\\ $x^2-3x=0$\\ $x(x-3)=0$\\ $x=0\, \text {hoặc}\, x=3$ \end {multicols} Vậy phương trình có 4 nghiệm là $x=-4,x=0,x=3,x=7$.  \hfill \qedEX 
\end{Solbook}
